\newcommand{\Q}{Q1} 

\newcounter{myWeekNum}
\stepcounter{myWeekNum}
%
\newcommand{\myWeek}{\themyWeekNum
    \stepcounter{myWeekNum}
    \ifnum\themyWeekNum=53
         \setcounter{myWeekNum}{1}
    \else\fi
}
\renewcommand{\baselinestretch}{0.9} 

\newcommand{\riskcell}[1]{%
    \FPeval\result{round(#1*100,0)} % Convert risk (0-1) to percentage for coloring
    \ifdim#1pt<0.3pt % If the risk is less than 0.3
        \cellcolor{green!30}#1%
    \else\ifdim#1pt<0.6pt % If the risk is between 0.3 and 0.6
        \cellcolor{yellow!30}#1%
    \else % If the risk is 0.6 or above
        \cellcolor{red!30}#1%
    \fi\fi
}
% \newcommand{\risk}[3]{#1 & \riskcell{#2} & \riskcell{#3}}

%
\crefname{deliverable}{deliverable}{deliverables}
\newtheorem{deliverable}{deliverable}

\def\mystartdate{2024-10-01}%starting date of the calendar
\def\myenddate{2025-05-30}%ending date of the calendar
\def\myxunit{0.7mm}%width of 1 day

\definecolor{gray}{HTML}{7F7F7F}
\definecolor{blue}{HTML}{1F77B4}

% \newcommand{\cc\1}{red!50} % phase 1 color
% \newcommand{\cc2}{magenta!50} % phase 2 color
% \newcommand{\cc3}{teal!50} % phase 3 color
% \newcommand{\cc4}{cyan!50} % phase 4 color
% \newcommand{\cc5}{blue!50} % phase 5 color
% \newcommand{\cc6}{purple!50} % phase 6 color
% \newcommand{\test123}{hi}

\usepackage{ifthen}  % For conditional statements

\newcommand{\Theme}[2]{%
  \ifthenelse{\equal{#1}{1}}{\textcolor{red!50}{#2}}{}%
      \ifthenelse{\equal{#1}{2}}{\textcolor{magenta!50}{#2}}{}%
      \ifthenelse{\equal{#1}{3}}{\textcolor{teal!50}{#2}}{}%
      \ifthenelse{\equal{#1}{4}}{\textcolor{cyan!50}{#2}}{}%
      \ifthenelse{\equal{#1}{5}}{\textcolor{blue!50}{#2}}{}%
      \ifthenelse{\equal{#1}{6}}{\textcolor{purple!50}{#2}}{}% 
}

\newcommand{\theme}[1]{%
      \ifthenelse{\equal{#1}{1}}{red!50}{}%
      \ifthenelse{\equal{#1}{2}}{magenta!50}{}%
      \ifthenelse{\equal{#1}{3}}{teal!50}{}%
      \ifthenelse{\equal{#1}{4}}{cyan!50}{}%
      \ifthenelse{\equal{#1}{5}}{blue!50}{}%
      \ifthenelse{\equal{#1}{6}}{purple!50}{}%
}


\newenvironment{conditions}[1][where:]
  {#1 \begin{tabular}[t]{>{$}l<{$} @{${}={}$} l}}
  {\end{tabular}\\[\belowdisplayskip]}

\makeatletter
% \renewcommand{\subsubsubsection}{\@startsection{paragraph}{4}{0ex}%
%    {-3.25ex plus -1ex minus -0.2ex}%
%    {1.5ex plus 0.2ex}%
%    {\normalfont\normalsize\bfseries}}
% \makeatother

\stepcounter{secnumdepth}
\stepcounter{tocdepth}

% \renewcommand{\today}{\the\year-\two@digits{\the\month}-\two@digits{\the\day}}
\renewcommand{\thedate}{\the\year-\two@digits{\the\month}-\two@digits{\the\day}}


\makeatletter
% Save the original \section command
\let\oldsection\section
% Redefine the \section command to handle both optional arguments and the star variant
\renewcommand{\section}{\@ifstar{\mysectionstar}{\@ifnextchar[{\mysectionopt}{\mysection}}}
% Define the version of \section with a star (unnumbered)
\newcommand{\mysectionstar}[1]{%
  \def\currentsec{#1}% Update to use a more specific variable name
  \oldsection*{#1}%
}
% Define the version of \section with an optional argument
\newcommand{\mysectionopt}[2][]{%
  \def\currentsec{#2}% Update to use a more specific variable name
  \oldsection[#1]{#2}%
}
% Define the version of \section without an optional argument
\newcommand{\mysection}[1]{%
  \def\currentsec{#1}% Update to use a more specific variable name
  \oldsection{#1}%
}

% Assuming similar functionality is desired for \subsection,
% you would follow a similar pattern:
\let\oldsubsection\subsection
\renewcommand{\subsection}{\@ifstar{\mysubsectionstar}{\@ifnextchar[{\mysubsectionopt}{\mysubsection}}}
\newcommand{\mysubsectionstar}[1]{%
  \def\currentsubsec{#1}% Update to use a more specific variable name
  \oldsubsection*{#1}%
}
\newcommand{\mysubsectionopt}[2][]{%
  \def\currentsubsec{#2}% Update to use a more specific variable name
  \oldsubsection[#1]{#2}%
}
\newcommand{\mysubsection}[1]{%
  \def\currentsubsec{#1}% Update to use a more specific variable name
  \oldsubsection{#1}%
}
\makeatother

\newcommand{\inputsubsection}[1]{%
  \subsection{#1}%
  \documentclass[11pt]{article}

\usepackage[colorlinks=true, allcolors=blue]{hyperref}
\usepackage[english]{babel}
\usepackage[numbers]{natbib}
\usepackage{url}
\usepackage[utf8x]{inputenc}
\usepackage{amsmath}
\usepackage{graphicx}
% \usepackage{parskip} 
\usepackage{fancyhdr}
\usepackage{vmargin}
\usepackage{tikz}
\usepackage{media9}
\usepackage{pgfplots} 
\usepackage{minted}
\usepackage{tikz}
\usepackage{pgfgantt}
\usepackage{pgfplots} 
\usepackage{titlesec}
\usepackage{caption}
\usepackage{subcaption}
 \usepackage[compatibility,siunitx,  americanvoltages, americancurrents, europeanresistors, europeaninductors, americanports,%
 straightlabels, fetbodydiode, straightvoltages]{circuitikz}
\usepackage{amssymb}
\usepackage{multirow}
\usepackage{siunitx}
\usepackage{pgfplots}
\usepackage{pgfgantt}
\usepackage{array}
\usepackage{longtable}
\usepackage{cleveref}
\usepackage{datetime}
\usepackage{tcolorbox}
\usepackage{physics}
\usepackage{amsmath}
\usepackage{tikz}
\usepackage{mathdots}
\usepackage{yhmath}
\usepackage{cancel}
\usepackage{color}
\usepackage{siunitx}
\usepackage{array}
\usepackage{multirow}
\usepackage{amssymb}
\usepackage{gensymb}
\usepackage{tabularx}
\usepackage{extarrows}
\usepackage{booktabs}
\usepackage{xcolor}
\usepackage{algorithm}
\usepackage{algpseudocode}
\usepackage{xpatch}
\usepackage{amsthm}
\usepackage{pdflscape}
\usepackage{rotating}
\usepackage[table]{xcolor} % For coloring cells
\usepackage{colortbl} 
\usepackage{fp}
% \usepackage[skip=0pt]{caption}
\usetikzlibrary{decorations.pathmorphing}
% \usetikzlibrary{decorations.pathreplacing}
\usetikzlibrary{fadings}
\usetikzlibrary{patterns}
\usetikzlibrary{shadows.blur}
\usetikzlibrary{shapes}
\usetikzlibrary{backgrounds}
\usetikzlibrary{positioning}

\usepackage{minted}
\usepackage{booktabs}
\usepackage{adjustbox} % For resizing tables and content
\usepackage{array} % For better column formatting


% \pgfplotsset{width=7cm,compat=1.3}
% \pgfplotsset{compat=newest}
% \pgfplotsset{plot coordinates/math parser=false} 
% \pgfplotsset{
%         % define the layers you need.
%         % (Don't forget to add `main' somewhere in that list!!)
%         layers/my layer set/.define layer set={
%             background,
%             main,
%             foreground
%         }{
%             % you could state styles here which should be moved to
%             % corresponding layers, but that is not necessary here.
%             % That is why we don't state anything here
%         },
%         % activate the newly created layer set
%         set layers=my layer set,
%     }

\input{Config/conf.tex}


\newcommand{\Q}{Q1} 

\newcounter{myWeekNum}
\stepcounter{myWeekNum}
%
\newcommand{\myWeek}{\themyWeekNum
    \stepcounter{myWeekNum}
    \ifnum\themyWeekNum=53
         \setcounter{myWeekNum}{1}
    \else\fi
}
\renewcommand{\baselinestretch}{0.9} 

\newcommand{\riskcell}[1]{%
    \FPeval\result{round(#1*100,0)} % Convert risk (0-1) to percentage for coloring
    \ifdim#1pt<0.3pt % If the risk is less than 0.3
        \cellcolor{green!30}#1%
    \else\ifdim#1pt<0.6pt % If the risk is between 0.3 and 0.6
        \cellcolor{yellow!30}#1%
    \else % If the risk is 0.6 or above
        \cellcolor{red!30}#1%
    \fi\fi
}
% \newcommand{\risk}[3]{#1 & \riskcell{#2} & \riskcell{#3}}

%
\crefname{deliverable}{deliverable}{deliverables}
\newtheorem{deliverable}{deliverable}

\def\mystartdate{2024-10-01}%starting date of the calendar
\def\myenddate{2025-05-30}%ending date of the calendar
\def\myxunit{0.7mm}%width of 1 day

\definecolor{gray}{HTML}{7F7F7F}
\definecolor{blue}{HTML}{1F77B4}

% \newcommand{\cc\1}{red!50} % phase 1 color
% \newcommand{\cc2}{magenta!50} % phase 2 color
% \newcommand{\cc3}{teal!50} % phase 3 color
% \newcommand{\cc4}{cyan!50} % phase 4 color
% \newcommand{\cc5}{blue!50} % phase 5 color
% \newcommand{\cc6}{purple!50} % phase 6 color
% \newcommand{\test123}{hi}

\usepackage{ifthen}  % For conditional statements

\newcommand{\Theme}[2]{%
  \ifthenelse{\equal{#1}{1}}{\textcolor{red!50}{#2}}{}%
      \ifthenelse{\equal{#1}{2}}{\textcolor{magenta!50}{#2}}{}%
      \ifthenelse{\equal{#1}{3}}{\textcolor{teal!50}{#2}}{}%
      \ifthenelse{\equal{#1}{4}}{\textcolor{cyan!50}{#2}}{}%
      \ifthenelse{\equal{#1}{5}}{\textcolor{blue!50}{#2}}{}%
      \ifthenelse{\equal{#1}{6}}{\textcolor{purple!50}{#2}}{}% 
}

\newcommand{\theme}[1]{%
      \ifthenelse{\equal{#1}{1}}{red!50}{}%
      \ifthenelse{\equal{#1}{2}}{magenta!50}{}%
      \ifthenelse{\equal{#1}{3}}{teal!50}{}%
      \ifthenelse{\equal{#1}{4}}{cyan!50}{}%
      \ifthenelse{\equal{#1}{5}}{blue!50}{}%
      \ifthenelse{\equal{#1}{6}}{purple!50}{}%
}


\newenvironment{conditions}[1][where:]
  {#1 \begin{tabular}[t]{>{$}l<{$} @{${}={}$} l}}
  {\end{tabular}\\[\belowdisplayskip]}

\makeatletter
% \renewcommand{\subsubsubsection}{\@startsection{paragraph}{4}{0ex}%
%    {-3.25ex plus -1ex minus -0.2ex}%
%    {1.5ex plus 0.2ex}%
%    {\normalfont\normalsize\bfseries}}
% \makeatother

\stepcounter{secnumdepth}
\stepcounter{tocdepth}

% \renewcommand{\today}{\the\year-\two@digits{\the\month}-\two@digits{\the\day}}
\renewcommand{\thedate}{\the\year-\two@digits{\the\month}-\two@digits{\the\day}}


\makeatletter
% Save the original \section command
\let\oldsection\section
% Redefine the \section command to handle both optional arguments and the star variant
\renewcommand{\section}{\@ifstar{\mysectionstar}{\@ifnextchar[{\mysectionopt}{\mysection}}}
% Define the version of \section with a star (unnumbered)
\newcommand{\mysectionstar}[1]{%
  \def\currentsec{#1}% Update to use a more specific variable name
  \oldsection*{#1}%
}
% Define the version of \section with an optional argument
\newcommand{\mysectionopt}[2][]{%
  \def\currentsec{#2}% Update to use a more specific variable name
  \oldsection[#1]{#2}%
}
% Define the version of \section without an optional argument
\newcommand{\mysection}[1]{%
  \def\currentsec{#1}% Update to use a more specific variable name
  \oldsection{#1}%
}

% Assuming similar functionality is desired for \subsection,
% you would follow a similar pattern:
\let\oldsubsection\subsection
\renewcommand{\subsection}{\@ifstar{\mysubsectionstar}{\@ifnextchar[{\mysubsectionopt}{\mysubsection}}}
\newcommand{\mysubsectionstar}[1]{%
  \def\currentsubsec{#1}% Update to use a more specific variable name
  \oldsubsection*{#1}%
}
\newcommand{\mysubsectionopt}[2][]{%
  \def\currentsubsec{#2}% Update to use a more specific variable name
  \oldsubsection[#1]{#2}%
}
\newcommand{\mysubsection}[1]{%
  \def\currentsubsec{#1}% Update to use a more specific variable name
  \oldsubsection{#1}%
}
\makeatother

\newcommand{\inputsubsection}[1]{%
  \subsection{#1}%
  \documentclass[11pt]{article}

\input{Config/inlcude.tex}
\input{Config/conf.tex}
\input{Config/macros.tex}
% \input{Config/macros.tex}
% Language setting
% Replace `english' with e.g. `spanish' to change the document language
\usepackage[english]{babel}

% Set page size and margins
% Replace `letterpaper' with `a4paper' for UK/EU standard size
% \usepackage[letterpaper,top=2cm,bottom=2cm,left=3cm,right=3cm,marginparwidth=1.75cm]{geometry}

% Useful packages
\usepackage{amsmath}
\usepackage{graphicx}
\usepackage[colorlinks=true, allcolors=blue]{hyperref}
\usepackage[cc]{titlepic}



\begin{document}

\input{Config/title_page.tex}
\pagebreak

\section{Outline of the Problem to be Solved}
\input{src/1.tex}


\section{Details of files used in the project}
\input{src/1.1.tex}


\section{Key Data Items}
\input{src/2.tex}

\section{Function Declarations}
\input{src/3.tex}

\section{Testing Information}
\input{src/4.tex}

\section{Flowcharts}
\input{src/5.tex}

\pagebreak
\end{document}
%
}

\newcommand{\inputsubsubsection}[1]{%
  \subsubsection{#1}%
  \input{Src/\currentsec/\currentsubsec/#1.tex}%
}

\newcommand{\Diagram}[2]{%
  \begin{figure}[h]
    \centering
    \includegraphics[width=0.8\textwidth]{Src/\currentsec/\currentsubsec/#1}
    \caption{#2}
    \label{fig:#1}
  \end{figure}
  \FloatBarrier
}


% 

\newcommand{\Q}{Q1} 

\newcounter{myWeekNum}
\stepcounter{myWeekNum}
%
\newcommand{\myWeek}{\themyWeekNum
    \stepcounter{myWeekNum}
    \ifnum\themyWeekNum=53
         \setcounter{myWeekNum}{1}
    \else\fi
}
\renewcommand{\baselinestretch}{0.9} 

\newcommand{\riskcell}[1]{%
    \FPeval\result{round(#1*100,0)} % Convert risk (0-1) to percentage for coloring
    \ifdim#1pt<0.3pt % If the risk is less than 0.3
        \cellcolor{green!30}#1%
    \else\ifdim#1pt<0.6pt % If the risk is between 0.3 and 0.6
        \cellcolor{yellow!30}#1%
    \else % If the risk is 0.6 or above
        \cellcolor{red!30}#1%
    \fi\fi
}
% \newcommand{\risk}[3]{#1 & \riskcell{#2} & \riskcell{#3}}

%
\crefname{deliverable}{deliverable}{deliverables}
\newtheorem{deliverable}{deliverable}

\def\mystartdate{2024-10-01}%starting date of the calendar
\def\myenddate{2025-05-30}%ending date of the calendar
\def\myxunit{0.7mm}%width of 1 day

\definecolor{gray}{HTML}{7F7F7F}
\definecolor{blue}{HTML}{1F77B4}

% \newcommand{\cc\1}{red!50} % phase 1 color
% \newcommand{\cc2}{magenta!50} % phase 2 color
% \newcommand{\cc3}{teal!50} % phase 3 color
% \newcommand{\cc4}{cyan!50} % phase 4 color
% \newcommand{\cc5}{blue!50} % phase 5 color
% \newcommand{\cc6}{purple!50} % phase 6 color
% \newcommand{\test123}{hi}

\usepackage{ifthen}  % For conditional statements

\newcommand{\Theme}[2]{%
  \ifthenelse{\equal{#1}{1}}{\textcolor{red!50}{#2}}{}%
      \ifthenelse{\equal{#1}{2}}{\textcolor{magenta!50}{#2}}{}%
      \ifthenelse{\equal{#1}{3}}{\textcolor{teal!50}{#2}}{}%
      \ifthenelse{\equal{#1}{4}}{\textcolor{cyan!50}{#2}}{}%
      \ifthenelse{\equal{#1}{5}}{\textcolor{blue!50}{#2}}{}%
      \ifthenelse{\equal{#1}{6}}{\textcolor{purple!50}{#2}}{}% 
}

\newcommand{\theme}[1]{%
      \ifthenelse{\equal{#1}{1}}{red!50}{}%
      \ifthenelse{\equal{#1}{2}}{magenta!50}{}%
      \ifthenelse{\equal{#1}{3}}{teal!50}{}%
      \ifthenelse{\equal{#1}{4}}{cyan!50}{}%
      \ifthenelse{\equal{#1}{5}}{blue!50}{}%
      \ifthenelse{\equal{#1}{6}}{purple!50}{}%
}


\newenvironment{conditions}[1][where:]
  {#1 \begin{tabular}[t]{>{$}l<{$} @{${}={}$} l}}
  {\end{tabular}\\[\belowdisplayskip]}

\makeatletter
% \renewcommand{\subsubsubsection}{\@startsection{paragraph}{4}{0ex}%
%    {-3.25ex plus -1ex minus -0.2ex}%
%    {1.5ex plus 0.2ex}%
%    {\normalfont\normalsize\bfseries}}
% \makeatother

\stepcounter{secnumdepth}
\stepcounter{tocdepth}

% \renewcommand{\today}{\the\year-\two@digits{\the\month}-\two@digits{\the\day}}
\renewcommand{\thedate}{\the\year-\two@digits{\the\month}-\two@digits{\the\day}}


\makeatletter
% Save the original \section command
\let\oldsection\section
% Redefine the \section command to handle both optional arguments and the star variant
\renewcommand{\section}{\@ifstar{\mysectionstar}{\@ifnextchar[{\mysectionopt}{\mysection}}}
% Define the version of \section with a star (unnumbered)
\newcommand{\mysectionstar}[1]{%
  \def\currentsec{#1}% Update to use a more specific variable name
  \oldsection*{#1}%
}
% Define the version of \section with an optional argument
\newcommand{\mysectionopt}[2][]{%
  \def\currentsec{#2}% Update to use a more specific variable name
  \oldsection[#1]{#2}%
}
% Define the version of \section without an optional argument
\newcommand{\mysection}[1]{%
  \def\currentsec{#1}% Update to use a more specific variable name
  \oldsection{#1}%
}

% Assuming similar functionality is desired for \subsection,
% you would follow a similar pattern:
\let\oldsubsection\subsection
\renewcommand{\subsection}{\@ifstar{\mysubsectionstar}{\@ifnextchar[{\mysubsectionopt}{\mysubsection}}}
\newcommand{\mysubsectionstar}[1]{%
  \def\currentsubsec{#1}% Update to use a more specific variable name
  \oldsubsection*{#1}%
}
\newcommand{\mysubsectionopt}[2][]{%
  \def\currentsubsec{#2}% Update to use a more specific variable name
  \oldsubsection[#1]{#2}%
}
\newcommand{\mysubsection}[1]{%
  \def\currentsubsec{#1}% Update to use a more specific variable name
  \oldsubsection{#1}%
}
\makeatother

\newcommand{\inputsubsection}[1]{%
  \subsection{#1}%
  \documentclass[11pt]{article}

\input{Config/inlcude.tex}
\input{Config/conf.tex}
\input{Config/macros.tex}
% \input{Config/macros.tex}
% Language setting
% Replace `english' with e.g. `spanish' to change the document language
\usepackage[english]{babel}

% Set page size and margins
% Replace `letterpaper' with `a4paper' for UK/EU standard size
% \usepackage[letterpaper,top=2cm,bottom=2cm,left=3cm,right=3cm,marginparwidth=1.75cm]{geometry}

% Useful packages
\usepackage{amsmath}
\usepackage{graphicx}
\usepackage[colorlinks=true, allcolors=blue]{hyperref}
\usepackage[cc]{titlepic}



\begin{document}

\input{Config/title_page.tex}
\pagebreak

\section{Outline of the Problem to be Solved}
\input{src/1.tex}


\section{Details of files used in the project}
\input{src/1.1.tex}


\section{Key Data Items}
\input{src/2.tex}

\section{Function Declarations}
\input{src/3.tex}

\section{Testing Information}
\input{src/4.tex}

\section{Flowcharts}
\input{src/5.tex}

\pagebreak
\end{document}
%
}

\newcommand{\inputsubsubsection}[1]{%
  \subsubsection{#1}%
  \input{Src/\currentsec/\currentsubsec/#1.tex}%
}

\newcommand{\Diagram}[2]{%
  \begin{figure}[h]
    \centering
    \includegraphics[width=0.8\textwidth]{Src/\currentsec/\currentsubsec/#1}
    \caption{#2}
    \label{fig:#1}
  \end{figure}
  \FloatBarrier
}


% Language setting
% Replace `english' with e.g. `spanish' to change the document language
\usepackage[english]{babel}

% Set page size and margins
% Replace `letterpaper' with `a4paper' for UK/EU standard size
% \usepackage[letterpaper,top=2cm,bottom=2cm,left=3cm,right=3cm,marginparwidth=1.75cm]{geometry}

% Useful packages
\usepackage{amsmath}
\usepackage{graphicx}
\usepackage[colorlinks=true, allcolors=blue]{hyperref}
\usepackage[cc]{titlepic}



\begin{document}

\begin{titlepage}
    \centering
    %\vspace*{0.5 cm}
    \includegraphics[scale = 0.4]{Config/uon.png}\\[2 cm]	% University Logo

    
    \textsc{\Large Department of Mechanical, Materials and Manufacturing Engineering}\\[0.3 cm]	% University Name
    \textsc{\large Computer Engineering and Mechatronics }\\[0.3 cm]				% Course Code
    \textsc{\large (MMME3085 UNUK) (AUT1 24-25)}\\[3 cm]				% Course Name

    \rule{\linewidth}{0.2 mm} \\[0.2 cm]
    { \huge \bfseries \thetitle}\\
    \rule{\linewidth}{0.2 mm} \\[3 cm]

    % \includegraphics[scale = 0.3]{Config/a.jpg}\\[1.0 cm]

    \large \emph{Author:} \theauthor \\[0.5 cm]


    {\large \today}\\[1 cm]

%     \textsc{\textbf{\large{Abbreviations:}}}\\[0.3 cm]
%     \normalsize
%     \textbf{UKF} - Unscented Kalman Filter, 
% \textbf{EKF} - Extended Kalman Filter, 
% \textbf{FMS} - Flight Management System, 
% \textbf{ECON} - economy, 
% \textbf{LIDAR} - Light Detection And Ranging, 
% \textbf{TEK} - Total Energy Kompensation, 
% \textbf{ASI} - Airspeed Indicator, 
% \textbf{AoA} - Angle of Attack, 
% \textbf{TRL} - Technology Readiness Level, 
% \textbf{NMEA} - National Marine Electronics Association, 
% \textbf{RASP} - Regional Atmospheric Soaring Prediction, 
% \textbf{GNSS} - Global Navigation Satellite System, 
% \textbf{GPS} - Global Positioning System, 
% \textbf{DMA} - Direct Memory Access,
\end{titlepage}

% \section{Overleaf Hints (to be deleted later)}

% \subsection{jumping between doc and latex}
% use ctr+click to jump in pages. \\

% use ctr+S to compile document.


% \subsection{Drawings}
% make drawings with: https://www.mathcha.io/editor/ \\


% \subsection{Graphs}
% all graphs should be done in matlab. Use matlab2tikz to create a latex file:

% https://uk.mathworks.com/matlabcentral/fileexchange/22022-matlab2tikz-matlab2tikz

% use the command to add your graph: 
% \begin{verbatim}
%         \input{path/to/your/file.tex}
% \end{verbatim}  

% This should be wrapped in a figure see \cref{sub:figures}

% \subsection{Tables}
% to make tables use: 
% https://www.tablesgenerator.com/

% use the [ht!] flag to specify table position.\\

% use following to ensure it stays in bounds
% \begin{verbatim}
%     \centering
%     \resizebox{\columnwidth}{!}
% \end{verbatim}

% if wraparound boxes are needed use:
% \begin{verbatim}
%     \tabular{p{100}} where p 100 is in pt.
% \end{verbatim}\\

% see \cref{tab:inital_spec_vs_Considerations} for example\\

% \subsection{Peer Review process}
% leave in green, for someone else to read. don't click accept on your own work \\

% if you see green that is not your own work. review it, then accept it. make sure to use a spell checker (you can get Grammarly for overleaf)\\

% \subsection{figures} \label{sub:figures}
% please use the following template for figures:

% \begin{figure}[ht!]
%     \centering
%     % \includegraphics{}
%     \caption{Caption}
%     \label{fig:enter-label}
% \end{figure}


% \subsection{referencing}
% \subsubsection{Cross referencing}

% use \verb|\Cref and \cref| ... eg \Cref{tab:inital_spec_vs_Considerations} and \cref{tab:inital_spec_vs_Considerations}

% label where appropriate:

% eg: \verb|\label{fig:figure_descritpion}| or \verb|\label{tab:table_description}| or \verb|\label{sec:section_name}|

% \subsubsection{citations}

% get extension BibItNow:

% https://addons.mozilla.org/en-US/firefox/addon/bibitnow/

% copy reference in bibtex format.

% paste in file biblist.bib found in the main file tree near bottom.

% eg:

% \begin{verbatim}
%     @misc{CanoeSlalom,
% 	title = {{Canoe Slalom}},
% 	journal = {ICF - Planet Canoe},
% 	year = {2023},
% 	month = oct,
% 	note = {[Online; accessed 31. Oct. 2023]},
% 	url = {https://www.canoeicf.com/disciplines/canoe-slalom}
% }
% \end{verbatim}


% Then use \verb|\cite{CanoeSlalom}| to reference: \cite{CanoeSlalom}

% if done correctly it should automatically appear hyper linked to the bottom of the document.



% \pagebreak
% \vspace*{\fill}
% \begin{abstract}
%    This report will discuss the feasibility of creating a full slalom canoe and kayak timing and scoring systems to automate the adjudication of the water sport. The report investigates the background and context behind the sport, including looking at existing solutions and analysing where similar automation has been achieved in other sporting disciplines. The stakeholder's requirements have been analysed qualitatively and quantitatively to generate a set of specifications for initial idea generation. An IoT radio localisation system has been conceptualised to map, track and report the athlete's position in time and space. Satellite nodes on each pole contain sensors to detect hits and act as radio beacons for localisation and data communication. An offline computer processes this data to generate scoring and timing data, displayed in real-time to a spectator through an app. Subsystems for rover and satellite prototypes, as well as a prototype scoring system and application, were implemented and tested. The prototypes were analysed against the specifications and original stakeholder requirements. Reflections were made on all aspects of the project and the project as a whole, including project management. Areas for future work and further development proposed. Broader context and commercialisation strategies were discussed to extrapolate the work's long-term implications. 
% \end{abstract}
% \vspace*{\fill}
% \pagebreak
\pagebreak

\section{Outline of the Problem to be Solved}
The program aims to print text on a robot arm by sending G-code instructions over a serial interface. It uses a font file that defines code for different ASCII characters. The program must ask the user for a text file to print and the text's height.

The project will be developed for a cross-platform experience using Windows (MSYS2), Linux, and Mac. GitHub will be used for version control. Automated workflows on Git Hub will compile and run the code to ensure platform compatibility. Documentation will be generated via doxygen in an automated workflow on Git Hub in the gh-pages branch, which will then be published on a static website. The program will be compiled using the GCC compiler, which makes it easy to compile on all platforms. 

The first aspect of the project will include parsing the font file. The program will use a structure to hold the code instructions. Ideally, this should use dynamic memory allocation to allow for different and more complex fonts with variable strokes/letters. The file will use the fopen() function from stdio.h to open the file. A loop will iterate through all file lines using the fgets() function. When 999 is found, indicating a new character, the program will allocate memory for the new character and the strokes that follow. All characters dynamically allocated will be mapped with a hash table, with the ASCII value of the character as the key.

The font file is defined by a series of 128 ASCII characters, which is normalised to a font size of 18mm. After the initial 999, followed by Ascii char and the number of strokes, the following lines define the pen's movements. The format is x position, y position, and pen up/down state. This must be converted to G code for the robot to operate with. For pen up, use "S1000 G1" and "S0 G0" for pen down, followed by the x and y positions. The robot must also be initiated with "G1 X0 Y0 F1000 $\backslash$n" to home the robot and set the feed rate, "M3$\backslash$n" is also used to start the spindle, used to enable the pen.

An interesting problem is selecting the G-code Information given the ASCII character. A simple approach would be to iterate through the array and find the matching ASCII character. This, however, would have a slow time complexity of O(n). A better approach might be to sort the array so the index of the code array is the same as the ASCII value. This would have a time complexity of O(1). However, this requires sorting the font file or implementing a sorting algorithm. Alternatively, a hash table could be used with an O(1) time complexity. This would allow nonsorted font files to be used efficiently and is easily extendable for other ASCII standards, including Unicode. 

The project's next stage will require scaling the code instructions; a simple command line dialog using printf() and scanf() will allow the user to input the font size they want. The program will then iterate through characters and font strokes and scale the x and y strokes. The scale is defined by double scale = (double)/ 18.0f; 

Next, the stage will be to input the test.txt file. A simple command line dialog using printf() and scanf() will allow users to input the file they want to print out. The file can be opened with fgets(). The specification requires that each word is read from the file individually and processed before reading the next one. While fscanf() could be used to read each word, it misses the ability to process punctuation or special characters such as a space or $\backslash$n. Therefore, the best way of processing will be to use fgetc() to read each character individually, appending it to a buffer. When a special character is found, the buffer is copied to an array, which can be processed. 

As each word is pulled from the text file, the program will call a function to assemble the code to send to the robot. Static variables will be utilized to keep track of the cursor position. The function will iterate through the word (array of characters), where it will first look for special characters, such as $\backslash$n, to update the cursor position(garage return with 5mm line spacing ). The ASCII value of the character will be used as the key in a hash function to return a pointer to the code instructions for the character. The function will then iterate through the stokes of the character, sending the code to the robot with the appropriate G0 or G1 command for the pen up or down position. Next, the cursor position will be updated to leave a small space between letters. Before sending a new G-code instruction, the program will wait for the OK$\backslash$n response from the robot to confirm it has room in the buffer for the next instruction. The robot will be communicated via a serial port using the RS232 library provided.




\section{Details of files used in the project}
\subsection{lib/rs232.c}

rs232 library written by Teunis van Beelen to enable rs232 communication with the robot. It has been modified to work with Mac. flags \_\_linux\_\_ , \_\_FreeBSD\_\_ , \_\_APPLE\_\_ and \_\_WINDOWS\_\_, set by the system are used at compile time to apply the correct implementation for the system. For Mac systems, RS232\_GetPortnr should be used to find the comp port number; else, manually define the comport lookup a table to include the correct port name, e.g., "/dev/tty.usbmodem1101".
\subsection{lib/rs232.h}
header file for rs232.c, contains the relevant header files for different systems included at compile time, and contains function prototypes for rs232.c.

\subsection{lib/serial.c}
Serial communication is a library communicating with the robot via the rs232 library. Compile time logic includes \#defines for Serial\_Mode, to be used when no external device is available to connect to. Also, \#DEBUG\_MODE is used to turn on and off debug print statements. If the system is Linux, this file will define Sleep() as a wrapper for the usleep () function available to Linux/Unix systems.
\subsection{lib/serial.h}
Header file for serial.c, contains function prototypes for serial.c and includes the relevant header files for the system at compile time. \#cport\_nr should be defined as the respective comport for Windows systems; for Mac and Linux, this value refers to the index of the comport in the port lookup table, which should be edited in the rs232.c file. additional \#bdrate\ defines the baud rate for the serial communication; for the robot, this should be 115200.

\subsection{src/main.c}
This is the main entry point for the program. This program initializes the robot, parses font data from a specified font file, and allows the user to specify a text height and input file containing the text to be drawn. It then processes the input text, converts it into G-code, and sends the commands to the robot to draw the text. Finally, it frees allocated resources and concludes the operation. User interface functions are also defined in this file.
\subsection{src/main.h}
header for main.c, contains function prototypes for user interface functions and includes the relevant header files for the program.

\subsection{src/robot/robot.c}
This file implements functions that control a robot through RS232 commands. It includes functionality to Initialize and start up the robot, Move the robot to a home position, Send stroke commands that move the pen or tool along specified coordinates, and toggle its state. The code is based on the example provided for the project.
\subsection{src/robot/robot.h}
Header file for robot.c, This file provides the interfaces for controlling the robot's position and movements. It includes functions to start up the robot, send strokes (pen movements) to it, and return it to a 'home' position. It also defines various constants that govern its coordinate space and text dimensions.

% #define MAX_X_VALUE_MM 100.0   /**< Maximum X coordinate value in millimeters. */
% #define MAX_Y_VALUE_MM 0       /**< Maximum Y coordinate value in millimeters. */
% #define MIN_X_VALUE_MM 0       /**< Minimum X coordinate value in millimeters. */
% #define MIN_Y_VALUE_MM -1000.0 /**< Minimum Y coordinate value in millimeters. */

% #define HOME_X_VALUE_MM 0.0 /**< Home position X coordinate in millimeters. */
% #define HOME_Y_VALUE_MM 0.0 /**< Home position Y coordinate in millimeters. */

% #define CHARACTER_SPACE_MM 18.0 /**< Default spacing in millimeters between characters. */
% #define LINE_SPACE_MM 5.0       /**< Default line spacing in millimeters. */

% #define MINIMUM_TEXT_HEIGHT_MM 4.0  /**< Minimum permissible text height in millimeters. */
% #define MAXIMUM_TEXT_HEIGHT_MM 10.0 /**< Maximum permissible text height in millimeters. */

% #define FONT_FILE "SingleStrokeFont.txt" /**< Default font file name. */

\begin{itemize}
    \setlength\itemsep{0pt}
    \setlength\parskip{0pt}
    \setlength\parsep{0pt}
    \item MAX\_X\_VALUE\_MM -- Maximum X coordinate value in millimeters.
    \item MAX\_Y\_VALUE\_MM -- Maximum Y coordinate value in millimeters.
    \item MIN\_X\_VALUE\_MM -- Minimum X coordinate value in millimeters.
    \item MIN\_Y\_VALUE\_MM -- Minimum Y coordinate value in millimeters.
    \item HOME\_X\_VALUE\_MM -- Home position X coordinate in millimeters.
    \item HOME\_Y\_VALUE\_MM -- Home position Y coordinate in millimeters.
    \item CHARACTER\_SPACE\_MM -- Default spacing in millimeters between characters.
    \item LINE\_SPACE\_MM -- Default line spacing in millimeters.
    \item MINIMUM\_TEXT\_HEIGHT\_MM -- Minimum permissible text height in millimeters.
    \item MAXIMUM\_TEXT\_HEIGHT\_MM -- Maximum permissible text height in millimeters.
    \item FONT\_FILE -- Default font file name.
\end{itemize}


\subsection{src/robot/gcode.c}
Implementation of text processing functions that generate G-code from text using font data. Processes a text string and generates corresponding G-code commands based on the provided font data. Reads a text file, processes each word, and utilizes `generate\_gcode` to convert the text into G-code.
\subsection{src/robot/gcode.h}
header file for gcode.c, contains declarations for text processing functions that utilize font data and robot movement.


\subsection{src/robot/cursor.c}
This file manages the cursor position and movement within predefined boundaries, providing functions to set, move, and return to a new line, handle carriage returns, update positioning, and check for word overflow conditions. It also contains the cursor object constructor, which is publicly accessible.
\subsection{src/robot/cursor.h}
Header file for cursor.c, contains declarations of the cursor object and constructor functions for cursor management.

\subsection{src/misc/coord.h}
Defines basic 2d coordinate type and basic operations for 2d coordinates, including addition, subtraction, multiplication, division, and scaling.
\subsection{src/misc/error.h}
Defines error codes for the program. It defines a function to handle errors and print messages to the console. requires \#PRINT\_ERROR\_CODES to be defined for printf statements to be compiled.

\subsection{src/font/fontChar.c}
This header file provides the declarations for the fontCharacter\_t structure and its associated functions. The fontCharacter\_t represents a font character defined by a sequence of strokes. Each stroke indicates a movement of the pen (with pen up or pen down states) to draw the character. The file also defines the fontCharConstructor function which creates and initializes a fontCharacter\_t instance.
\subsection{src/font/fontChar.h}
This file provides the functionality to store and manage font characters as a series of strokes. It includes functions to append and free character strokes and a constructor to create a new font character.

\subsection{src/font/dataData.c}
This file provides the functionality to store and manage font characters inside a hash table. Each font character is associated with an ASCII key and is stored as a series of strokes. The hash table uses chaining to handle collisions, allowing multiple font characters with the same hash to be stored as a linked list of nodes. Functions are provided to insert and look up font characters, scale the strokes of all font characters by a given factor, parse a font file to populate the hash table, and free all allocated resources.
\subsection{src/font/fontData.h}
This header defines the fontData\_t structure and related components for storing and managing font characters in a hash table. Each entry in the hash table corresponds to an ASCII character and its associated fontCharacter\_t object. The fontData\_t structure provides function pointers for operations such as inserting, looking up, scaling, parsing from a file, and freeing the entire font data structure. The size of the ASCII standard is defined as 128 in a macro.


\section{Key Data Items}

% \begin{minted}[bgcolor=cyan!5,fontsize=\small]{c}
% typedef struct stroke_s{
%     double x;
%     double y;
%     bool pen_state;
% } stroke_t;

% typedef struct fontCharacter_s{
%     char ascii_key;
%     uint8_t num_strokes;
%     stroke_t *strokes;
% } fontCharacter_t;

% typedef struct cursor_s {
%     double x, y;
% }cursor_t;

% typedef struct hashNode_s {
%     char key;
%     fontCharacter_t *character;
%     struct hashNode_s *next;
% } hashNode_t;

% typedef struct fontData_s {
%     double scale:
%     hashNode_t *table[128]; // Hash table
% } fontData_t;
% \end{minted}

\begin{minted}[bgcolor=cyan!5,fontsize=\small,breaklines]{c}


typedef struct Coord2D_s
{
    double x;
    double y;
} Coord2D_t;

typedef Coord2D_t Vect2d_t;

typedef struct stroke_s
{
    Vect2d_t vec;   
    bool pen_state; 
} stroke_t;

typedef struct hashNode_s
{
    char key;                   
    fontCharacter_t *character; 
    struct hashNode_s *next;   
} hashNode_t;

typedef struct fontData_s
{
    double fontScale;                    
    hashNode_t *table[ASCII_CHARACTERS]; 

    errorCode_t (*free)(struct fontData_s *self);
    const fontCharacter_t *(*lookup)(const struct fontData_s *const self, const char ascii_key);
    errorCode_t (*scale)(struct fontData_s *const self, double scale);
    errorCode_t (*insert)(struct fontData_s *const self, fontCharacter_t *const character);
    errorCode_t (*parse)(struct fontData_s *const self, const char *filename);
} fontData_t;



typedef struct fontCharacter_s
{
    char asciiKey;      
    uint8_t numStrokes; 
    uint8_t strokeIdx;  
    stroke_t *strokes;  

    errorCode_t (*appendStroke)(struct fontCharacter_s *const self, const stroke_t stroke);
    errorCode_t (*free)(struct fontCharacter_s *self);
} fontCharacter_t;


  typedef struct cursor_s
{
    Coord2D_t posisiton;    
    Coord2D_t homePosition; 
    Coord2D_t maxPosition;  
    Coord2D_t minPosition;  
    bool init;              
    double scale;           
    double lineSpace;       
    double characterSpace;  

    errorCode_t (*set)(struct cursor_s *const self, const Coord2D_t position);
    errorCode_t (*move)(struct cursor_s *const self, const Coord2D_t delta);
    errorCode_t (*newline)(struct cursor_s *const self);
    errorCode_t (*carriagereturn)(struct cursor_s *const self);
    errorCode_t (*update)(struct cursor_s *const self);
    bool (*testWordOverflow)(const struct cursor_s *const self, const char *const word);
} cursor_t;

typedef enum errorCode_e
{
    SUCCESS = 0,                    /**< Operation completed successfully. */
    ERROR_OPEN_FILE,                /**< Error opening text file. */
    ERROR_INVALID_SCALE_INPUT,      /**< Invalid scale input (out of allowed range). */
    ERROR_NO_TEXT_FILE,             /**< No text file found. */
    ERROR_NO_FONT_DATA,             /**< No font data available. */
    CURSOR_OUT_OF_BOUNDS,           /**< Cursor position is outside the valid range. */
    FONT_CHARACTER_NOT_FOUND,       /**< Required font character not found in data. */
    ERROR_MEMORY_ALLOCATION_FAILED, /**< Memory allocation failure. */
    ERROR_NULL_POINTER,             /**< Null pointer encountered. */
    ERROR_UNABLE_TO_OPEN_COM_PORT,  /**< Unable to open the COM port. */
    ERROR_OUT_OF_BOUNDS,            /**< Out-of-bounds error (generic). */
    WORD_TOO_LONG,                  /**< Provided word is too long. */
    ERROR_INVALID_INPUT,            /**< Invalid input provided. */
    ERROR_INVALID_FILE,             /**< Invalid file. */
    ERROR_INVALID_FONT_FILE,        /**< Invalid font file. */
    ERROR_INVALID_FONT_CHARACTER,   /**< Invalid font character definition. */
    ERROR_INVALID_FONT_STROKE,      /**< Invalid font stroke definition. */
    ERROR_INVALID_FONT_STROKE_VEC,  /**< Invalid font stroke vector. */
    ERROR_INSERT_CHARACTER,         /**< Error inserting character into data structure. */
    ERROR_APPEND_STROKE,            /**< Error appending stroke to character. */
    ERROR_PARSE_STROKE,             /**< Error parsing stroke definition. */
    ERROR_UNEXPECTED_EOF,           /**< Unexpected end-of-file encountered. */
    ERROR_PARSE_CHARACTER           /**< Error parsing character definition. */
} errorCode_t;
    \end{minted}
\begin{table}[H]
    \centering
    \adjustbox{max width=\textwidth}{%
    \begin{tabular}{|>{\raggedright\arraybackslash}m{0.15\textwidth}|>{\raggedright\arraybackslash}m{0.16\textwidth}|>{\raggedright\arraybackslash}m{0.75\textwidth}|}
    \hline
    \textbf{Name}         & \textbf{Data Type}               & \textbf{Rationale}                                                                                      \\ \hline

    \textbf{file} & 
    FILE
    & part of stdio.h, used to define a file which can be opened with fopen() and worked on with other functions\\ \hline

    \textbf{position, homePosition, maxPosition, minPosition}& 
    Coord2D\_t& 2d coordinate system, with double presision to ensure maximum operational precision\\ \hline 

 \textbf{vec}& Vect2d\_t&used to form the stroke\_t. inherites Coord2D\_t to allow for polymorphism with functions that take Coord2D\_t as their argument.\\\hline


    \textbf{strokes}       & 
    stroke\_t& Structure structure to hold a single stroke, x and y are double to preserve high presision, and the pen state is bolean because it has 2 states, on or off. \\ \hline
    
    \textbf{character}& 
    fontCharacter\_t*
    & Holds the array of strokes\_t and the metadata for a single character. Because the strokes are dynamically allocated memory, the number of strokes is required to know the length of the array of strokes. A char is one byte and, therefore, can hold all combinations of 128 ASCII codes. num\_strokes is assigned as an 8-bit unsigned integer, where negative values are impossible, and more than 255 is unlikely. memory will be allocated using malloc, so the data type is a pointer hence acts as array. Additionaly function pointers, for abstraction and polymorphism, are used to call static functions. overloading of free() as well as an append() function.\\ \hline
    
    \textbf{Cursor}       & 
    static cursor\_t
    &  Tracks the current position of the "pen" on the robot's surface for precise placement of letters. Also manages the pen, including limiting bounds and applying special characters such as newline, and carridgereturn. Coord2D\_t allows for floating point presision to define, positions, and robot boundarys.. double presision is used to define scale, line space and character space which is precomputed at compile time. will be declared as static within the generate\_gcode function to keep track of the cursor position, thefore a boolean member init, is declared for lasy initilsation. functions pointers include: set, move, newline, carriagereturn and update. These are used to access static functions through the cursor object.\\ \hline

    \textbf{hashTable} & 
    hashNode\_t*
    & a pointer to the hash table array, contains the key which is used in the hash function to find the correct FontCharacter pointer. although not neccassary becasue the data ascii font file data is sequencial, a HashNode\_s structure pointer is used for collision handling. memory will be allocated using malloc, so the data type is a pointer hence acts as array.\\ \hline
    
    \textbf{FontData}& 
    fontData\_t*
    & A container for holding the hash table array. also holds the scale factor to be used for determining white space. double to preserve presision. function pointers: free, lookup, scale, insert and parse are also used to access static functions where the class is implemented.\\ \hline
 \textbf{error}& errorCode\_t & enum because it is a list of grouped identitys where the value is meaningless. contains all of the error codes for the program. is used by ErrorHandler() to display errors when they occour. \\ \hline
    \end{tabular}
    }
    \caption{Primary Data Types for the G-code Printing Program}
    \label{tab:data_types}
    \end{table}



\section{Function Declarations}

    

    \begin{minted}[bgcolor=cyan!5,fontsize=\small]{c}
int hash_table_insert(fontData_t *hash_table, const char ascii_key, fontCharacter_t *character);
        \end{minted}
        \textbf{Breif:} Inserts a font character into the hash table. \\
        \textbf{Parameters:} \\
        \hspace*{1cm}\texttt{hash\_table} -- Pointer to the hash table. \\
        \hspace*{1cm}\texttt{ascii\_key} -- ASCII key for the font character. \\
        \hspace*{1cm}\texttt{character} -- Pointer to the font character to insert. \\
        \textbf{Returns:} \\
        \hspace*{1cm}0 on success, -1 on failure. \\[1em]
        
        \begin{minted}[bgcolor=cyan!5,fontsize=\small]{c}
const fontCharacter_t *hash_table_lookup(const fontData_t *hash_table, const char ascii_key);
        \end{minted}
        \textbf{Breif:} Looks up a font character in the hash table by ASCII key. \\
        \textbf{Parameters:} \\
        \hspace*{1cm}\texttt{hash\_table} -- Pointer to the hash table. \\
        \hspace*{1cm}\texttt{ascii\_key} -- ASCII key of the font character to retrieve. \\
        \textbf{Returns:} \\
        \hspace*{1cm}Pointer to the corresponding fontCharacter\_t, or \texttt{NULL} if not found. \\[1em]
        
        \begin{minted}[bgcolor=cyan!5,fontsize=\small]{c}
void parse_font_file(FontData *font_data, const char *filename);
        \end{minted}
        \textbf{Breif:} Parses the font file and populates a hash table with font characters. \\
        \textbf{Parameters:} \\
        \hspace*{1cm}\texttt{font\_data} -- Pointer to the font data to populate. \\
        \hspace*{1cm}\texttt{filename} -- Name of the font file to parse. \\
        
        \begin{minted}[bgcolor=cyan!5,fontsize=\small]{c}
void scale_font_data(FontData *font_data);
        \end{minted}
        \textbf{Breif:} Scales the font data by a given factor. \\
        \textbf{Parameters:} \\
        \hspace*{1cm}\texttt{font\_data} -- Pointer to the font data to scale. \\

        
        \begin{minted}[bgcolor=cyan!5,fontsize=\small]{c}
void process_text_file(const FontData *font_data);
        \end{minted}
        \textbf{Breif:} Processes a text file and generates G-code for the text. \\
        \textbf{Parameters:} \\
        \hspace*{1cm}\texttt{font\_data} -- Pointer to the font data to use, also containing scale for whitespace between characters. \\[1em]
        
        \begin{minted}[bgcolor=cyan!5,fontsize=\small]{c}
void generate_gcode(const FontData *font_data, const char *text);
        \end{minted}
        \textbf{Breif:} Generates G-code for a given text string. \\
        \textbf{Parameters:} \\
        \hspace*{1cm}\texttt{font\_data} -- Pointer to the font data to use, also containing scale for whitespace between characters. \\
        \hspace*{1cm}\texttt{text} -- Text string (single word) to generate G-code for.

        \begin{minted}[bgcolor=cyan!5,fontsize=\small]{c}
void free_font_data(fontData_t *font_data);
        \end{minted}
        \textbf{Breif:} Frees the memory allocated for the font data. itterate through all hashNode\_t,fontCharacter\_t and strokes\_t to free memory allocated by malloc\\
        \textbf{Parameters:} \\
        \hspace*{1cm}\texttt{font\_data} -- Pointer to the font data to be freed.

\section{Testing Information}

\begin{longtable}{|>{\raggedright\arraybackslash}m{0.17\textwidth}|>{\raggedright\arraybackslash}m{0.2\textwidth}|>{\raggedright\arraybackslash}m{0.2\textwidth}|>{\raggedright\arraybackslash}m{0.35\textwidth}|}
    \hline
    \textbf{Function} & \textbf{Test Case} & \textbf{Test Data} & \textbf{Expected Output} \\ \hline
    parse\_font\_file & empty filename & filename="" & Error message: "Filename is empty" \\ \hline
    parse\_font\_file & fopen() fails & filename="font.exe" & Error message: "unable to open file" \\ \hline
    parse\_font\_file & memory allocation fails & Simulate malloc failure  &  Error message: "Memory allocation failed" \\ \hline
    parse\_font\_file & normal case & filename= "SingleStrokeFont.txt" &  normal operation: calls hash\_table\_insert(), 128 times \\ \hline
    hash\_table\_insert & NULL hastable & hash\_table=NULL & Error message: "Hash table is NULL" \\ \hline
    hash\_table\_insert & NULL character & character=NULL & Error message: "Character is NULL" \\ \hline
    hash\_table\_insert & memory allocation fails & Simulate malloc failure  &  Error message: "Memory allocation failed" \\ \hline
    hash\_table\_insert & normal case & ascii\_key=0, character=fontCharacter\_t &  normal operation: inserts fontCharacter into hash table \\ \hline
    hash\_table\_insert & insert multiple characters & insert 'A', 'B', 'C' &  normal operation: inserts fontCharacter into hash table \\ \hline
    hash\_table\_insert & confirm correct operation with hash\_table\_lookup & *font\_char = hash\_table\_lookup( font\_data, key) & font\_char$->$ascii\_key == Key \\ \hline
    hash\_table\_lookup & NULL hastable & hash\_table=NULL & Error message: "Hash table is NULL" \\ \hline
    hash\_table\_lookup & normal case & ascii\_key=0 &  normal operation: returns fontCharacter from hash table with key 0\\ \hline
    hash\_table\_lookup & look up character that doesnt exist & ascii\_key=999 &  Error message: "Character not found" \\ \hline
    scale\_font\_data & NULL font\_data & font\_data=NULL & Error message: "Font data is NULL" \\ \hline
    scale\_font\_data & input not between 4-10mm & 11 & Error message: "Input not between 4-10mm" \\ \hline
    scale\_font\_data & test with 'A' 4mm & \footnotesize$\begin{array}{ccc}
        999 &65 &6 \\
        0 &0 &0\\
        6 &18& 1\\
        12 &0& 1\\
        3 &9& 0\\
        9& 9& 1\\
        18 &0& 0\\
        \end{array}$ & \footnotesize$\begin{array}{ccc}
        999 &65 &6 \\
        0 &0 &0\\
        1.33 &4.00& 1\\
        2.67 &0& 1\\
        0.67 &2.00& 0\\
        2.00& 2.00& 1\\
        4.00 &0& 0\\
        \end{array}$ \\ \hline
        scale\_font\_data & test with 'A' 10mm & \footnotesize$\begin{array}{ccc}
            999 &65 &6 \\
            0 &0 &0\\
            6 &18& 1\\
            12 &0& 1\\
            3 &9& 0\\
            9& 9& 1\\
            18 &0& 0\\
            \end{array}$ & \footnotesize$\begin{array}{ccc}
            999 &65 &6 \\
            0 &0 &0\\
            3.33 &10.00& 1\\
            6.67 &0& 1\\
            1.67 &5.00& 0\\
            5.00& 5.00& 1\\
            10.00 &0& 0\\
            \end{array}$ \\ \hline
    process\_text\_file & NULL font\_data & font\_data=NULL & Error message: "Font data is NULL" \\ \hline
    process\_text\_file & fopen() fails & filename="text.exe" & Error message: "unable to open file" \\ \hline
    process\_text\_file & Test each word read indvidually & "Hello World" &  normal operation: calls generate\_gcode() twice \\ \hline
    generate\_gcode & NULL font\_data & font\_data=NULL & Error message: "Font data is NULL" \\ \hline
    generate\_gcode & test $\backslash$n mechanism & text[0] = "$\backslash$n" & cursor moves to next line with 5mm line spaace\\ \hline
    generate\_gcode & test normal case & text[0] = "A" &  normal operation: calls hash\_table\_lookup() and sends G-code to robot, waits of ok response after every send. \\ \hline
    generate\_gcode & 4mm test 'A' output & text[0] = "A " & \footnotesize$\begin{array}{cccl}
        G0 &X0.00 &Y0.00; & pen\ up\\
        G1 &X1.33 &Y4.00; & pen\ down\\
        G1 &X2.67 &Y0.00; & pen\ down\\
        G0 &X0.67 &Y2.00; & pen\ up\\
        G1 &X2.00 &Y2.00; & pen\ down\\
        G0 &X4.00 &Y0.00; & pen\ down\\
        G0 &x0 &y0; & pen\ up\\
        \end{array}$ \\ \hline
    generate\_gcode & 10mm test 'A' output & text[0] = "A " & \footnotesize$\begin{array}{cccl}
        G0 &X0.00 &Y0.00; & pen\ up\\
        G1 &X3.33 &Y10.00; & pen\ down\\
        G1 &X6.67 &Y0.00; & pen\ down\\
        G0 &X1.67 &Y5.00; & pen\ up\\
        G1 &X5.00 &Y5.00; & pen\ down\\
        G0 &X10.00 &Y0.00; & pen\ down\\
        G0 &x0 &y0; & pen\ up\\
        \end{array}$ \\ \hline
    generate\_gcode & test multiple letters checking for proper character spacing with 4mm spacing & text[0] = "AA " &  \footnotesize$\begin{array}{cccl}
        G0 &X0.00 &Y0.00; & pen\ up\\
        G1 &X1.33 &Y4.00; & pen\ down\\
        G1 &X2.67 &Y0.00; & pen\ down\\
        G0 &X0.67 &Y2.00; & pen\ up\\
        G1 &X2.00 &Y2.00; & pen\ down\\
        G0 &X4.00 &Y0.00; & pen\ down\\
        G0 &X4.00 &Y0.00; & pen\ up\\
        G1 &X5.33 &Y4.00; & pen\ down\\
        G1 &X6.67 &Y0.00; & pen\ down\\
        G0 &X4.67 &Y2.00; & pen\ up\\
        G1 &X6.00 &Y2.00; & pen\ down\\
        G0 &X8.00 &Y0.00; & pen\ up\\
        G0 &X12.00 &Y0.00; & pen\ up\\
        G0 &X0 &Y0; & pen\ up\\
        \end{array}$ \\ \hline



    main & test normal case & SingleStrokeFont, test.txt: "Hello World", 4mm text height & \scriptsize$\begin{array}{cccl}
    G0 &X0.00 Y&0.00  &Pen up\\
    G1 &X0.00 Y&4.00  &Pen down\\
    G0 &X2.67 Y&0.00  &Pen up\\
    G1 &X2.67 Y&4.00  &Pen down\\
    G0 &X0.00 Y&2.00  &Pen up\\
    G1 &X2.67 Y&2.00  &Pen down\\
    G0 &X4.00 Y&0.00  &Pen up\\
    G0 &X4.00 Y&1.33  &Pen up\\
    G1 &X6.67 Y&1.56  &Pen down\\
    G1 &X6.00 Y&2.67  &Pen down\\
    G1 &X4.67 Y&2.67  &Pen down\\
    G1 &X4.00 Y&2.00  &Pen down\\
    G1 &X4.00 Y&0.44  &Pen down\\
    G1 &X4.67 Y&0.00  &Pen down\\
    G1 &X6.00 Y&0.00  &Pen down\\
    G1 &X6.67 Y&0.44  &Pen down\\
    G0 &X8.00 Y&0.00  &Pen up\\
    G0 &X8.67 Y&0.00  &Pen up\\
    G1 &X10.00 &Y0.00; &Pen down\\
    G0 &X9.33 Y&0.00  &Pen up\\
    G1 &X9.33 Y&4.00  &Pen down\\
    G1 &X8.67 Y&4.00  &Pen down\\
    G0 &X12.00 &Y0.00; &Pen up\\
    G0 &X12.67 &Y0.00; &Pen up\\
    G1 &X14.00 &Y0.00; &Pen down\\
    G0 &X13.33 &Y0.00; &Pen up\\
    G1 &X13.33 &Y4.00; &Pen down\\
    G1 &X12.67 &Y4.00; &Pen down\\
    G0 &X16.00 &Y0.00; &Pen up\\
    G0 &X17.33 &Y0.00; &Pen up\\
    G1 &X16.00 &Y0.44; &Pen down\\
    G1 &X16.00 &Y2.00; &Pen down\\
    G1 &X17.33 &Y2.44; &Pen down\\
    G1 &X18.67 &Y2.00; &Pen down\\
    G1 &X18.67 &Y0.44; &Pen down\\
    G1 &X17.33 &Y0.00; &Pen down\\
    G0 &X20.00 &Y0.00; &Pen up\\
    G0 &X24.00 &Y0.00; &Pen up\\
    G0 &X24.00 &Y4.00; &Pen up\\
    G1 &X24.67 &Y0.00; &Pen down\\
    G1 &X25.33 &Y3.11; &Pen down\\
    G1 &X26.00 &Y0.00; &Pen down\\
    G1 &X26.67 &Y4.00; &Pen down\\
    G0 &X28.00 &Y0.00; &Pen up\\
    G0 &X29.33 &Y0.00; &Pen up\\
    G1 &X28.00 &Y0.44; &Pen down\\
    G1 &X28.00 &Y2.00; &Pen down\\
    G1 &X29.33 &Y2.44; &Pen down\\
    G1 &X30.67 &Y2.00; &Pen down\\
    G1 &X30.67 &Y0.44; &Pen down\\
    G1 &X29.33 &Y0.00; &Pen down\\
    G0 &X32.00 &Y0.00; &Pen up\\
    G0 &X32.00 &Y0.00; &Pen up\\
    G1 &X32.00 &Y2.44; &Pen down\\
    G0 &X32.00 &Y1.78; &Pen up\\
    G1 &X33.33 &Y2.44; &Pen down\\
    G1 &X34.67 &Y1.78; &Pen down\\
    G0 &X36.00 &Y0.00; &Pen up\\
    G0 &X36.67 &Y0.00; &Pen up\\
    G1 &X38.00 &Y0.00; &Pen down\\
    G0 &X37.33 &Y0.00; &Pen up\\
    G1 &X37.33 &Y4.00; &Pen down\\
    G1 &X36.67 &Y4.00; &Pen down\\
    G0 &X40.00 &Y0.00; &Pen up\\
    G0 &X42.67 &Y0.44; &Pen up\\
    G1 &X41.33 &Y0.00; &Pen down\\
    G1 &X40.00 &Y0.44; &Pen down\\
    G1 &X40.00 &Y2.00; &Pen down\\
    G1 &X41.33 &Y2.44; &Pen down\\
    G1 &X42.67 &Y2.00; &Pen down\\
    G0 &X42.67 &Y4.00; &Pen up\\
    G1 &X42.67 &Y0.00; &Pen down\\
    G0 &X44.00 &Y0.00; &Pen up\\
    G0 &X44.00 &Y0.00; &Pen up\\
    G0 &X0 &Y0 ;& to origin\\
    \end{array}$\\ \hline
    main & test normal case with carrage return & SingleStrokeFont, test.txt: "Hello $\backslash$n World", 10mm text height & \scriptsize$\begin{array}{cccl}
        G0 &X0.00  &Y0.00;   &Pen up\\
        G1 &X0.00  &Y10.00;  &Pen down\\
        G0 &X6.67  &Y0.00;   &Pen up\\
        G1 &X6.67  &Y10.00;  &Pen down\\
        G0 &X0.00  &Y5.00;   &Pen up\\
        G1 &X6.67  &Y5.00;   &Pen down\\
        G0 &X10.00 &Y0.00;   &Pen up\\
        G0 &X10.00 &Y3.33;   &Pen up\\
        G1 &X16.67 &Y3.89;   &Pen down\\
        G1 &X15.00 &Y6.67;   &Pen down\\
        G1 &X11.67 &Y6.67;   &Pen down\\
        G1 &X10.00 &Y5.00;   &Pen down\\
        G1 &X10.00 &Y1.11;   &Pen down\\
        G1 &X11.67 &Y0.00;   &Pen down\\
        G1 &X15.00 &Y0.00;   &Pen down\\
        G1 &X16.67 &Y1.11;   &Pen down\\
        G0 &X20.00 &Y0.00;   &Pen up\\
        G0 &X21.67 &Y0.00;   &Pen up\\
        G1 &X25.00 &Y0.00;   &Pen down\\
        G0 &X23.33 &Y0.00;   &Pen up\\
        G1 &X23.33 &Y10.00;  &Pen down\\
        G1 &X21.67 &Y10.00;  &Pen down\\
        G0 &X30.00 &Y0.00;   &Pen up\\
        G0 &X31.67 &Y0.00;   &Pen up\\
        G1 &X35.00 &Y0.00;   &Pen down\\
        G0 &X33.33 &Y0.00;   &Pen up\\
        G1 &X33.33 &Y10.00;  &Pen down\\
        G1 &X31.67 &Y10.00;  &Pen down\\
        G0 &X40.00 &Y0.00;   &Pen up\\
        G0 &X43.33 &Y0.00;   &Pen up\\
        G1 &X40.00 &Y1.11;   &Pen down\\
        G1 &X40.00 &Y5.00;   &Pen down\\
        G1 &X43.33 &Y6.11;   &Pen down\\
        G1 &X46.67 &Y5.00;   &Pen down\\
        G1 &X46.67 &Y1.11;   &Pen down\\
        G1 &X43.33 &Y0.00;   &Pen down\\
        G0 &X50.00 &Y0.00;   &Pen up\\
        G0 &X60.00 &Y0.00;   &Pen up\\
        G0 &X60.00 &Y0.00;   &Pen up\\
        G0 &X0.00  &Y-5.00;  &Pen up\\
        G1 &X1.67  &Y-15.00; &Pen down\\
        G1 &X3.33  &Y-7.22;  &Pen down\\
        G1 &X5.00  &Y-15.00; &Pen down\\
        G1 &X6.67  &Y-5.00;  &Pen down\\
        G0 &X10.00 &Y-15.00; &Pen up\\
        G0 &X13.33 &Y-15.00; &Pen up\\
        G1 &X10.00 &Y-13.89; &Pen down\\
        G1 &X10.00 &Y-10.00; &Pen down\\
        G1 &X13.33 &Y-8.89;  &Pen down\\
        G1 &X16.67 &Y-10.00; &Pen down\\
        G1 &X16.67 &Y-13.89; &Pen down\\
        G1 &X13.33 &Y-15.00; &Pen down\\
        G0 &X20.00 &Y-15.00; &Pen up\\
        G0 &X20.00 &Y-15.00; &Pen up\\
        G1 &X20.00 &Y-8.89;  &Pen down\\
        G0 &X20.00 &Y-10.56; &Pen up\\
        G1 &X23.33 &Y-8.89;  &Pen down\\
        G1 &X26.67 &Y-10.56; &Pen down\\
        G0 &X30.00 &Y-15.00; &Pen up\\
        G0 &X31.67 &Y-15.00; &Pen up\\
        G1 &X35.00 &Y-15.00; &Pen down\\
        G0 &X33.33 &Y-15.00; &Pen up\\
        G1 &X33.33 &Y-5.00;  &Pen down\\
        G1 &X31.67 &Y-5.00;  &Pen down\\
        G0 &X40.00 &Y-15.00; &Pen up\\
        G0 &X46.67 &Y-13.89; &Pen up\\
        G1 &X43.33 &Y-15.00; &Pen down\\
        G1 &X40.00 &Y-13.89; &Pen down\\
        G1 &X40.00 &Y-10.00; &Pen down\\
        G1 &X43.33 &Y-8.89;  &Pen down\\
        G1 &X46.67 &Y-10.00; &Pen down\\
        G0 &X46.67 &Y-5.00;  &Pen up\\
        G1 &X46.67 &Y-15.00; &Pen down\\
        G0 &X50.00 &Y-15.00; &Pen up\\
        G0 &X50.00 &Y-15.00; &Pen up\\
        G0 &X0 &Y0; & toorigin \\
    \end{array}$\\ \hline
        \end{longtable}



\section{Flowcharts}
please find flow charts in folder named RobotFlowchartGRD\_ID.pdf


\pagebreak
\end{document}
%
}

\newcommand{\inputsubsubsection}[1]{%
  \subsubsection{#1}%
  \input{Src/\currentsec/\currentsubsec/#1.tex}%
}

\newcommand{\Diagram}[2]{%
  \begin{figure}[h]
    \centering
    \includegraphics[width=0.8\textwidth]{Src/\currentsec/\currentsubsec/#1}
    \caption{#2}
    \label{fig:#1}
  \end{figure}
  \FloatBarrier
}

